\documentclass[11pt, a4paper]{article}

\usepackage[utf8]{inputenc}
\usepackage{amsmath, amssymb, amsfonts}
\usepackage{geometry}
\geometry{margin=1in}
\usepackage{parskip}
\usepackage{color}

\title{\textbf{Problem Set 5: Lagos-Wright Model}}
\author{}
\date{}

\begin{document}

\maketitle

\section*{Exercise 1: Design of Central Bank Digital Currencies}

\subsection*{Part 1: Bellman Equations}

\subsubsection*{1. CM problem}

In the CM, a buyer with portfolio $(z_m,z_d)$ chooses consumption $x$, labor $h$, and next period's portfolio $(z'_m,z'_d)$.  
The Bellman equation is
\[
W(z_m,z_d)
= \max_{x,h,z'_m,z'_d} \{ x - h + \beta V(z'_m,z'_d) \}
\]
subject to the CM budget constraint
\[
x + z'_m + z'_d = h + z_m + (1+i_d) z_d + T.
\]

\subsubsection*{2. Linearity of the CM value function}

Use the budget constraint to eliminate labor:
\[
h = x + z'_m + z'_d - z_m - (1+i_d)z_d - T.
\]
Substituting into the Bellman equation gives
\begin{align*}
W(z_m,z_d)
&= \max_{x,z'_m,z'_d} \bigl\{
x - [x + z'_m + z'_d - z_m - (1+i_d)z_d - T]
+ \beta V(z'_m,z'_d) \bigr\} \\
&= \max_{x,z'_m,z'_d} \bigl\{
z_m + (1+i_d)z_d + T - z'_m - z'_d + \beta V(z'_m,z'_d)
\bigr\}.
\end{align*}

Define the constant
\[
W(0,0) \equiv \max_{x,z'_m,z'_d}
\{ T - z'_m - z'_d + \beta V(z'_m,z'_d) \},
\]
which does not depend on $(z_m,z_d)$.  
Then we can write
\[
W(z_m,z_d) = z_m + (1+i_d)z_d + W(0,0),
\]
showing that $W$ is linear in the current portfolio.

\subsubsection*{3. DM problem and its simplification}

In the DM, the buyer enters with balances $(z'_m,z'_d)$.  
With probabilities $\alpha_1,\alpha_2,\alpha_3$ they meet a seller of type 1, 2, or 3 respectively; with probability $1-\alpha_1-\alpha_2-\alpha_3$ there is no trade.  
In a type-$j$ meeting the buyer trades $q_j$ units, obtaining utility $u(q_j)$ and paying real balances $q_j$ in the instrument accepted at that meeting.

The DM Bellman equation is
\begin{align*}
V(z'_m,z'_d)
&= \alpha_1\bigl[u(q_1) + W(z'_m - q_1, z'_d)\bigr]
+ \alpha_2\bigl[u(q_2) + W(z'_m, z'_d - q_2)\bigr] \\
&\quad + \alpha_3\bigl[u(q_3) + W(z'_m - q_3, z'_d)\bigr]
+ (1-\alpha_1-\alpha_2-\alpha_3) W(z'_m,z'_d).
\end{align*}

Using linearity of $W$ (e.g.\ $W(z'_m - q_1,z'_d)=W(z'_m,z'_d)-q_1$) we get
\begin{align*}
V(z'_m,z'_d)
&= \alpha_1\bigl[u(q_1) + W(z'_m,z'_d) - q_1\bigr]
+ \alpha_2\bigl[u(q_2) + W(z'_m,z'_d) - q_2\bigr] \\
&\quad + \alpha_3\bigl[u(q_3) + W(z'_m,z'_d) - q_3\bigr]
+ (1-\alpha_1-\alpha_2-\alpha_3) W(z'_m,z'_d) \\
&= W(z'_m,z'_d)
+ \alpha_1[u(q_1)-q_1]
+ \alpha_2[u(q_2)-q_2]
+ \alpha_3[u(q_3)-q_3].
\end{align*}

Thus the DM value equals the CM value at the start-of-DM portfolio plus the expected net surplus from DM trade:
\[
V(z'_m,z'_d)
= W(z'_m,z'_d)
+ \alpha_1[u(q_1)-q_1]
+ \alpha_2[u(q_2)-q_2]
+ \alpha_3[u(q_3)-q_3].
\]

\subsubsection*{4. Deriving the CM objective in terms of liquidity}

Plug this expression into the CM problem:
\begin{align*}
W(z_m,z_d)
&= \max_{x,z'_m,z'_d}
\bigl\{ z_m + (1+i_d)z_d + T - z'_m - z'_d \\
&\qquad + \beta\bigl[z'_m + (1+i_d)z'_d + W(0,0)\bigr] \\
&\qquad + \beta\bigl(\alpha_1[u(q_1)-q_1]
+ \alpha_2[u(q_2)-q_2] + \alpha_3[u(q_3)-q_3]\bigr)
\bigr\}.
\end{align*}

Using the Fisher equation $1+i = 1/\beta$ (zero inflation) we can rewrite the coefficients on $z'_m,z'_d$ as proportional to $-i$ and $-(i-i_d)$ respectively, so up to irrelevant positive factors the buyer chooses $(z'_m,z'_d)$ to maximize
\[
\alpha_1[u(q_1)-q_1]
+ \alpha_2[u(q_2)-q_2]
+ \alpha_3[u(q_3)-q_3]
- i z'_m - (i - i_d) z'_d.
\]

Thus, maximizing the CM value is equivalent to maximizing expected DM surplus net of the opportunity cost of holding liquidity.

\subsection*{Part 2: Baseline Economy (No CBDC)}

In the baseline, there is only physical cash $z_m$ and deposits $z_d$.  
We assume buyers bring exactly enough cash to cover type-1 trades and deposits to cover type-2 and type-3 trades:
\[
z_m = q_1, \qquad z_d = \max\{q_2,q_3\},
\]
and in type-3 meetings they use the cheapest instrument first.  
Since $i > i - i_d$, deposits are cheaper than cash, so they strictly prefer to pay with deposits in type-3 meetings.

\subsubsection*{1. First-order conditions}

With these assumptions, the relevant objective over $(z_m,z_d)$ is
\[
\alpha_1[u(z_m)-z_m]
+ \alpha_2[u(z_d)-z_d]
+ \alpha_3[u(z_d)-z_d]
- i z_m - (i - i_d) z_d.
\]

The FOC with respect to $z_m$ is
\[
\alpha_1[u'(z_m)-1] - i = 0
\quad\Longrightarrow\quad
u'(z_m) = \frac{1+i}{\alpha_1}.
\]

The FOC with respect to $z_d$ is
\[
(\alpha_2+\alpha_3)[u'(z_d)-1] - (i - i_d) = 0
\quad\Longrightarrow\quad
u'(z_d) = \frac{1 + i - i_d}{\alpha_2+\alpha_3}.
\]

These are the money-demand conditions stated in the problem.

\subsubsection*{2. Effect of an increase in $i_d$ on $z_d$}

From the deposit FOC,
\[
u'(z_d) = \frac{1 + i - i_d}{\alpha_2+\alpha_3}.
\]
For fixed $i,\alpha_2,\alpha_3$, a higher $i_d$ lowers the right-hand side.  
Because $u''(q)<0$, $u'(q)$ is strictly decreasing in $q$, so a lower $u'(z_d)$ implies a higher $z_d$.  
Therefore, an increase in the private bank interest rate $i_d$ raises the aggregate demand for bank deposits.

\subsection*{Part 3: Introducing a Deposit-like CBDC}

Now the central bank introduces a deposit-like CBDC balance $z_c$ paying interest $i_c$.  
It uses the same digital infrastructure as bank deposits, so it is accepted in type-2 and type-3 meetings but not in type-1 meetings.  
Suppose $i_c > i_d$, so CBDC has a lower opportunity cost $i-i_c$ than private deposits $i-i_d$.

\subsubsection*{1. Objective with $(z_m,z_d,z_c)$}

With three instruments, the opportunity costs per unit are:
\[
\text{cash: } i, \qquad
\text{deposits: } i - i_d, \qquad
\text{CBDC: } i - i_c.
\]

The CM objective (up to constants) becomes
\[
\max_{z_m,z_d,z_c}
\left\{
\alpha_1[u(q_1)-q_1]
+ \alpha_2[u(q_2)-q_2]
+ \alpha_3[u(q_3)-q_3]
- i z_m - (i - i_d) z_d - (i - i_c) z_c
\right\}.
\]

\subsubsection*{2. Crowding out of bank deposits when $i_c>i_d$}

Suppose $i_c>i_d$, so that the CBDC has a lower opportunity cost than private deposits: $i-i_c<i-i_d$.

If both deposits and CBDC are used in type‑2 and type‑3 meetings (an interior solution), their marginal utilities–net‑of‑cost must be equal. Analogous to Part 2, the FOCs take the form

$$
u'(z_d)=\frac{1+i-i_d}{\alpha_2+\alpha_3},\quad
u'(z_c)=\frac{1+i-i_c}{\alpha_2+\alpha_3}.
$$

For an interior optimum in which both $z_d$ and $z_c$ are positive and used in the same types of meetings, we must have $u'(z_d)=u'(z_c)$. This requires

$$
\frac{1+i-i_d}{\alpha_2+\alpha_3}=\frac{1+i-i_c}{\alpha_2+\alpha_3}
\quad\Rightarrow\quad i_d=i_c,
$$

which contradicts the assumption $i_c>i_d$.

Therefore both assets cannot be held simultaneously in an interior solution: one must be at a corner. Since $i_c>i_d$ makes CBDC strictly cheaper (lower opportunity cost) than deposits, the buyer optimally holds only CBDC for type‑2 and type‑3 payments, implying

$$
z_d=0,\; z_c>0.
$$

Economically, the deposit‑like CBDC substitutes for bank deposits as a liquid, interest‑bearing means of payment in the DM. Buyers move their liquidity from private bank deposits to CBDC, shrinking the deposit base of banks—this is the “disintermediation” or “crowding‑out” risk of a deposit‑like CBDC design.

\subsection*{Part 4: Introducing a Cash-like CBDC}

Now suppose the central bank instead introduces a \emph{cash-like} CBDC.  
This digital currency does not require a bank account, is peer-to-peer, and is accepted in type-1 and type-3 meetings.  
It pays an interest rate $i_c$ with $0 < i_c < i_d$.  
Physical cash pays zero interest and has opportunity cost $i$, so a unit of cash-like CBDC has opportunity cost $i - i_c < i$ and strictly dominates physical cash in the meetings where both can be used.

\subsubsection*{1. New FOC for cash-like CBDC}

Because CBDC dominates physical cash, buyers optimally set $z_m = 0$ and use cash-like CBDC balances $z_c$ for type-1 meetings (and in type-3 meetings alongside deposits).  
The analogue of the cash FOC is
\[
\alpha_1[u'(z_c) - 1] - (i - i_c) = 0,
\]
which implies
\[
u'(z_c) = \frac{1 + i - i_c}{\alpha_1}.
\]

This matches the FOC requested in the problem statement.

\subsubsection*{2. Payment efficiency: showing $z_c > z_m$}

In the baseline without CBDC, the cash FOC was
\[
u'(z_m) = \frac{1+i}{\alpha_1}.
\]

With cash-like CBDC, the FOC is
\[
u'(z_c) = \frac{1 + i - i_c}{\alpha_1}.
\]

Because $0 < i_c < i_d$, we have $1 + i - i_c < 1 + i$, so
\[
u'(z_c) < u'(z_m).
\]

Since $u''(q)<0$, the marginal utility $u'(q)$ is strictly decreasing in $q$, so a lower marginal utility corresponds to a higher quantity:
\[
z_c > z_m.
\]

Thus agents hold strictly more real balances in the form of cash-like CBDC than they held as physical cash, allowing larger and more efficient DM trades in type-1 (and partly type-3) meetings; this is an increase in payment efficiency.

\subsubsection*{3. General equilibrium “crowding in” of deposits (intuition)}

The cash-like CBDC competes with physical cash but not perfectly with deposits, because deposits are accepted in type-2 and type-3 meetings, while cash-like CBDC is accepted in type-1 and type-3 meetings.  
By lowering the opportunity cost of liquidity in type-1 and type-3 meetings, the CBDC raises the volume and efficiency of DM trade, increasing overall economic activity.  
Higher total trade and income increase the demand for payment services that can only be provided with deposits in type-2 meetings, which tends to raise the desired stock of bank deposits $z_d$.  
In general equilibrium, a well-designed cash-like CBDC can therefore produce a \emph{crowding-in} effect on the banking sector: despite competing with cash, it may increase, rather than decrease, the demand for private bank deposits.


\section*{Exercise 2: Token Pricing -- Speculation vs. Usage}

\subsection*{1. Token Asset Pricing Equation}

Start from Euler equation \((4)\) and divide both sides by \(p_t\):
\[
1
= \beta \frac{E_t[p_{t+1}]}{p_t}
\left(
1 + \alpha\left[\frac{u'(q)}{w'(q)} - 1\right]
\right).
\]

Using \(\beta = \frac{1}{1+i}\) from \((1)\) and \(\frac{E_t[p_{t+1}]}{p_t} = 1 + \pi^e\), we obtain
\[
1
= \frac{1}{1+i} (1+\pi^e)
\left(
1 + \alpha\left[\frac{u'(q)}{w'(q)} - 1\right]
\right).
\]

Multiply both sides by \(\frac{1+i}{1+\pi^e}\):
\[
\frac{1+i}{1+\pi^e}
=
1 + \alpha\left[\frac{u'(q)}{w'(q)} - 1\right].
\]

Hence the required liquidity premium is
\[
\alpha\left[\frac{u'(q)}{w'(q)} - 1\right]
= \frac{1+i}{1+\pi^e} - 1.
\]

This expression involves only the dollar interest rate \(i\) and the expected token appreciation \(\pi^e\).

\paragraph{Economic interpretation of \(\frac{1+i}{1+\pi^e}\).}

The term \(1+i\) is the gross return from holding dollars (the risk-free nominal asset), while \(1+\pi^e\) is the expected gross return from holding tokens as a financial asset.  
The ratio \(\frac{1+i}{1+\pi^e}\) measures how attractive dollars are relative to tokens on pure return grounds.  
If this ratio exceeds 1, holding dollars dominates financially, so the token must offer a positive liquidity premium (transaction/use value) to be held in equilibrium.

\subsection*{2. Pure Speculation (\(\alpha = 0\))}

Suppose the platform currently provides no utility from services, so \(\alpha = 0\).  
Then the liquidity premium is zero:
\[
\alpha\left[\frac{u'(q)}{w'(q)} - 1\right] = 0.
\]

Plugging into the asset-pricing relation yields
\[
0 = \frac{1+i}{1+\pi^e} - 1
\quad\Longrightarrow\quad
\frac{1+i}{1+\pi^e} = 1
\quad\Longrightarrow\quad
1+i = 1+\pi^e
\quad\Longrightarrow\quad
\pi^e = i.
\]

So if \(\alpha=0\) and users still willingly hold tokens in the CM, we must have
\[
\pi^e = i.
\]

\paragraph{Economic interpretation.}

With \(\alpha=0\), tokens have no current use as a medium of exchange: they deliver no service utility.  
Investors are indifferent between dollars and tokens only if the expected capital gain on tokens equals the dollar interest rate.  
In this case the token functions purely as a speculative asset, not as a medium of exchange.

\subsection*{3. Volatility and Velocity (\(\pi^e < 0\))}

Now assume the token is highly volatile and the market expects its price to crash, so \(\pi^e < 0\).  

From the pricing equation:
\[
\alpha\left[\frac{u'(q)}{w'(q)} - 1\right]
= \frac{1+i}{1+\pi^e} - 1.
\]

When \(\pi^e < 0\), we have \(1 + \pi^e < 1\), so
\[
\frac{1+i}{1+\pi^e} > 1+i > 1,
\]
and therefore the liquidity premium on the right-hand side is strictly positive and large:
\[
\alpha\left[\frac{u'(q)}{w'(q)} - 1\right] > 0 \quad \text{and large if } \pi^e \text{ is very negative.}
\]

Given \(\alpha>0\), this implies
\[
\frac{u'(q)}{w'(q)} - 1
= \frac{1}{\alpha}
\left(
\frac{1+i}{1+\pi^e} - 1
\right)
\]
must be large, so \(\frac{u'(q)}{w'(q)}\) must be high.  

For a standard concave \(u(\cdot)\), a high \(u'(q)\) corresponds to low \(q\): each unit of service has very high marginal utility.  
Thus to justify holding an asset with negative expected return, the per-unit liquidity benefit from spending it must be very large.

\paragraph{Economic intuition: velocity of circulation.}

A negative expected return \(\pi^e<0\) makes tokens unattractive as a store of value.  
Rational users hold only the minimal token inventory necessary to execute imminent transactions; they buy tokens just before using the service and spend them immediately.  
This behavior implies a high velocity of circulation: tokens are held for very short periods and used strictly as a transaction device, not for saving.

\subsection*{4. The Tipping Point for the CIA Constraint}

The cash-in-advance constraint
\[
w(q_{t+1}) \le p_{t+1} m_t
\]
binds when users hold just enough tokens to pay for services.  
The constraint ceases to strictly bind when the marginal value of liquidity is zero, i.e.\ when the liquidity premium vanishes:
\[
\alpha\left[\frac{u'(q)}{w'(q)} - 1\right] = 0.
\]

With \(\alpha>0\), this is equivalent to \(\frac{u'(q)}{w'(q)} = 1\), and from the asset-pricing relation
\[
\alpha\left[\frac{u'(q)}{w'(q)} - 1\right]
= \frac{1+i}{1+\pi^e} - 1 = 0
\]
we obtain the tipping point
\[
\frac{1+i}{1+\pi^e} = 1
\quad\Longrightarrow\quad
1+i = 1+\pi^e
\quad\Longrightarrow\quad
\pi^e = i.
\]

Thus, the cash-in-advance constraint ceases to strictly bind exactly when
\[
\boxed{\pi^e = i.}
\]

\paragraph{Economic interpretation.}

At \(\pi^e = i\), the token's expected capital gain equals the dollar interest rate.  
Even if the CIA constraint is slack, investors are willing to hold tokens purely as a financial asset, since they are as attractive as dollars.  

When \(\pi^e > i\), tokens are strictly more attractive than dollars as an investment:

\begin{itemize}
\item Agents hold tokens not only for transactions but also for speculative motives.
\item They accumulate token balances beyond what is needed to satisfy \(w(q_{t+1})\), so the CIA constraint becomes slack.
\item The token increasingly behaves as a pure speculative asset: users hoard more tokens than needed for service consumption, and velocity falls as balances are held for longer periods.
\end{itemize}


\end{document}